
\noindent 

Non-negative Matrix Factorization (NMF) is an unsupervised learning technique that factorizes a non-negative data matrix into the product of two smaller non-negative matrices. In the context of image analysis, such as face recognition, NMF is valued for producing a parts-based representation of the data – for example, decomposing facial images into intuitive building blocks like eyes, nose, and mouth \cite{diaz2021analysis}. Despite its usefulness, a key challenge in deploying NMF on real-world data is the presence of noise and corruptions. In practical scenarios, measurements and images are often contaminated by various types of noise, which can significantly degrade the quality of the factorization. NMF, particularly in its basic form, tends to be sensitive to large errors in the data due to the nature of its standard loss function \cite{kang2023robustness}. As a result, evaluating and improving the robustness of NMF under noisy conditions is critical for its reliability in applications like face image analysis.

In this study, we investigate the robustness of NMF algorithms when face image data are corrupted with salt-and-pepper noise, a common form of impulse noise. Salt-and-pepper noise randomly flips a fraction of pixels to extreme values (black or white), simulating pepper-like and salt-like speckles on the image. We focus on two well-known face datasets – the ORL face database and the Extended Yale B face database – as testbeds for our analysis. We systematically introduce noise at different severity levels, controlled by $p$ and $r$, simulating conditions from mild to severe \cite{kang2023robustness}. 

We compare two NMF approaches in this work: (1) the standard $L_2$-norm NMF, and (2) an $L_{2,1}$-norm based NMF algorithm, which is a robust variant that employs an $L_{2,1}$-norm loss function to reduce the influence of outliers and noise \cite{kong2011robust,lu2016projective}. We hypothesize that this robust formulation will produce more stable factorizations under high noise, maintaining better reconstruction accuracy than the traditional approach \cite{wu2018manifold}.